% ch17-systems.tex -- Low-latency design patterns
\chapter{Low-latency design patterns}
\label{ch:systems17}
\index{low latency}\index{ring buffer}\index{jitter}

Low latency isn't a mean problem --- it's a jitter problem. The
patterns in this chapter each remove a source of \emph{occasional}
slowness: a lock that sometimes convoys, an allocator that
sometimes walks its free list, a branch that sometimes misses.
The chapter runs the standard claims on real hardware first,
because two of the four come out backwards.

\section{The folklore audit}
\begin{figure}[htbp]
\centering
\begin{tabular}{lll}
\toprule
low-latency claim & this machine & verdict \\
\midrule
``branchless beats branchy'' & 25.1 vs 23.3\,ms & false (predictable) \\
``binary beats linear'' & 0.005 vs 0.020\,\textmu s (n=64) & true, $4\times$ \\
``malloc is slow'' & 11\,ns/pair & mostly false (t-cache) \\
``pool reuse wins'' & 0.000 vs 0.011\,\textmu s & true, margin small \\
\bottomrule
\end{tabular}
\caption{The folklore audit, ch17 edition: four claims run on real
hardware. ``Branchless'' loses on predictable branches; malloc's
thread cache makes the allocation scare-claim mostly myth ---
but predictable-branch cost and the small-size search result are
the parts everyone gets backwards.}
\label{fig:audit17}
\end{figure}


\begin{example}[Branchless loses on predictable data]
\texttt{ch17\_folklore.c}: a branchy \texttt{min} over predictable
inputs runs 23.3\,ms; the bit-trick branchless version runs
25.1\,ms --- 7\% \emph{slower}, because the predictor made the
branch free and the ALU work isn't. Branchless only pays where
the data defeats the predictor; the claim fails by inversion.
\end{example}

\begin{figure}[htbp]
\centering
\begin{tabular}{lcc}
\toprule
data pattern & branchy min & branchless min \\
\midrule
predictable (our run) & 23.3\,ms & 25.1\,ms \\
random (would need) & slow & faster \\
\bottomrule
\end{tabular}
\caption{Branchless only wins when the predictor loses: our data was
predictable, so the branch was free and the extra ALU work cost
7\%. The claim fails by inversion --- measure the data's
randomness first.}
\label{fig:branch17}
\end{figure}


\section{Search at both scales}
\begin{figure}[htbp]
\centering
\begin{tabular}{lrr}
\toprule
search & n=64 & n=16M \\
\midrule
linear & 0.020\,\textmu s & 2.0\,ms \\
binary & 0.005\,\textmu s & 0.10\,\textmu s \\
\bottomrule
\end{tabular}
\caption{Binary wins even at n=64 --- but by $4\times$, not the
20\,000$\times$ it wins at 16M. The ``linear is fine for small n''
folklore is true in ratio, false in direction: binary never lost.}
\label{fig:search17}
\end{figure}


Binary search wins at n=64 too --- but $4\times$, not the
20\,000$\times$ it wins at 16M. ``Linear is fine for small n'' is
true in ratio and false in direction; the honest version is
``linear's disadvantage is small at small n.''

\section{The allocation myth, priced}
Malloc on a thread-cached allocator costs 11\,ns/pair --- the
``malloc is slow'' scare is mostly myth at this granularity. But
pool reuse still wins (0.011 vs 0.000\,\textmu s) and, more
importantly, removes the allocator's tail: the reason low-latency
code pre-allocates isn't steady-state cost, it's the rare
free-list walk the p99 sees.

\section{The SPSC ring}
\begin{figure}[htbp]
\centering
\begin{tabular}{l|ccc}
\toprule
slot & 0 & \dots & 255 \\
\midrule
producer writes & head & & \\
consumer reads & & tail & \\
no lock needed & \multicolumn{3}{c}{one writer, one reader} \\
\bottomrule
\end{tabular}
\caption{The SPSC ring: power-of-two slots, head written only by the
producer, tail only by the consumer. No mutex, no allocation ---
the workhorse of real-time audio, networking, and our measured
1M-pair run.}
\label{fig:ring17}
\end{figure}


Power-of-two slots, head written only by the producer, tail only
by the consumer --- no mutex because each index has one owner.
The real \texttt{ch17\_ringbuf.c}: 1M push/pop pairs, 0
mismatches. It's the workhorse of audio pipelines and network
drivers precisely because it trades capacity headroom (a fixed
256 slots) for zero synchronisation.

\section{The cost sheet}
\begin{figure}[htbp]
\centering
\begin{tabular}{ll}
\toprule
what low latency costs & where you pay it \\
\midrule
jitter budget & pre-allocated pools, no syscalls \\
capacity headroom & fixed-size rings drop or block \\
code complexity & branchless/lock-free is harder \\
measurement rigor & p99 not mean --- ch15's lesson \\
\bottomrule
\end{tabular}
\caption{The real cost sheet: low-latency work is paid in headroom
and complexity, not in clever instructions --- and audited with
the tail-latency tools of ch15.}
\label{fig:cost17}
\end{figure}


\section{Common mistakes}
\begin{itemize}
\item \textbf{Branchless everywhere.} Our branchless \texttt{min}
      lost 7\% on predictable data --- the trick pays only where
      the predictor fails.
\item \textbf{``Never malloc in the hot path'' as dogma.} The
      thread cache is 11\,ns; the real reason is the allocator's
      \emph{tail}, not its mean.
\item \textbf{Lock-free = fast, assumed.} ch13 already priced it:
      contended lines cost more than locks at scale.
\item \textbf{Optimising the mean.} The user on the p99 feels the
      rare syscall/alloc/miss --- mean optimisation is often
      invisible.
\item \textbf{Unbounded queues for backpressure.} A ring that
      blocks at 256 beats a queue that swells to 2\,GB and then
      falls over --- capacity is a design decision.
\end{itemize}

\section{Summary}
\begin{tabularx}{\linewidth}{@{}l X@{}}
\toprule
\textbf{Piece} & \textbf{Verified version} \\
\midrule
jitter   & the real target --- p99 not mean \\
branchless & loses 7\% on predictable data (real) \\
search   & binary wins even at n=64 ($4\times$) \\
malloc   & 11\,ns/pair --- myth-busted, pool still wins \\
ring     & 1M pairs, 0 mismatches, no lock \\
\bottomrule
\end{tabularx}

\section*{Further reading}
\begin{itemize}
\item Mechanical Sympathy blog (jitter systems): \\ \texttt{mechanical-sympathy.blogspot.com}
\item OSTEP concurrency chapters (SPSC groundwork): \texttt{pages.cs.wisc.edu/\~{}remzi}
\end{itemize}

\begin{rulebox}[Watch instead]
\begin{itemize}
\item \textbf{Low Latency Engineering --- CppCon talk} --- the
      ring-buffer and jitter material at C++ depth
      (verified 2026-10-11).
      \url{https://youtu.be/NH1Tta7purM}
\end{itemize}
\end{rulebox}

\section*{Notes}
\begin{itemize}
\item All numbers verbatim \texttt{ch17\_folklore.c} on Apple
      Silicon; the branchless loss and malloc-myth rows are the
      chapter's honest surprises.
\item Ring buffer verified single-threaded; SPSC correctness under
      real concurrency needs the memory-ordering argument of the
      [hard] exercise.
\end{itemize}

\section*{Exercises}
\addcontentsline{toc}{section}{Exercises}
\stepcounter{section}

\begin{exercise}[easy]
Why is jitter, not the mean, the low-latency target?
\end{exercise}
\begin{solution}[hint]
The p99 is what the user sees; mean cost is usually fine already.
\end{solution}
\begin{exercise}[easy]
Our branchless min lost on predictable data --- which side of
the hardware made the branchy version free?
\end{exercise}
\begin{solution}[hint]
The branch predictor; it only loses on unpredictable data.
\end{solution}
\begin{exercise}[easy]
What makes the SPSC ring lock-free with only two fields?
\end{exercise}
\begin{solution}[hint]
Single owner per index --- no shared write to order.
\end{solution}
\begin{exercise}[easy]
Name the real cost the ring buffer pays for zero locks.
\end{exercise}
\begin{solution}[hint]
Fixed capacity --- it drops or blocks at 256 entries.
\end{solution}
\begin{exercise}[medium]
Binary beat linear at n=64 by only $4\times$. Explain why the
ratio is small there and huge at 16M, in memory terms.
\end{exercise}
\begin{solution}[hint]
n=64 fits cache/branch-predictor; at 16M each linear probe is a
cache miss vs binary's 24.
\end{solution}
\begin{exercise}[medium]
A teammate removes all mallocs from a hot loop and sees no
improvement. Using our 11\,ns figure, argue why --- and where the
win actually lives.
\end{exercise}
\begin{solution}[hint]
At 11\,ns each the mean cost is noise; the win is the tail ---
free-list walks and lock contention at the p99.
\end{solution}
\begin{exercise}[medium]
Extend the ring to MPMC (multi-producer/consumer). Which fields
stop being single-owner, and what primitive replaces them?
\end{exercise}
\begin{solution}[hint]
Head and tail become contended --- CAS loops or sequence numbers
(LMAX-style) restore order.
\end{solution}
\begin{exercise}[medium]
Design the experiment that shows WHERE branchless actually wins
on your machine. What data property flips the result?
\end{exercise}
\begin{solution}[hint]
Random vs sorted inputs --- mispredicts at $\sim$50\% each cost
10--20 cycles; predictable costs 0.
\end{solution}
\begin{exercise}[hard]
Argue the memory-ordering contract that makes our SPSC ring
correct on a weakly-ordered ISA: what must the producer's
head-store publish, and what can go wrong if the buffer write
isn't ordered before it? Give the Dekker-style reasoning.
\end{exercise}
\begin{solution}[hint]
The data write must be release-ordered before head; otherwise
the consumer reads tail-but-stale buffer (a torn read).
\end{solution}
\begin{exercise}[hard]
The folklore audit killed 2 of 4 claims. Write the claim about
low-latency design that IS true in your experience, run it, and
report the number --- and explain what makes it survive where
the others didn't.
\end{exercise}
\begin{solution}[hint]
True claims specify \emph{conditions} (``branchless wins on
random data'') --- they survive because the measurement matches
the assumption.
\end{solution}
\begin{exercise}[easy]
What does ``allocation-free'' actually mean in the ring code?
\end{exercise}
\begin{solution}[hint]
Fixed compile-time storage --- the 256-slot array lives in the
struct, not the heap.
\end{solution}
\begin{exercise}[medium]
Backpressure: our push returns -1 when full. Compare that policy
with blocking and with unbounded growth --- which is correct for
real-time audio and why?
\end{exercise}
\begin{solution}[hint]
Drop/overwrite --- a late audio sample is a click; a growing
queue is unbounded latency.
\end{solution}
\begin{exercise}[medium]
The ring uses \texttt{volatile} head/tail. Why isn't that enough
on a real concurrent ISA --- what's missing?
\end{exercise}
\begin{solution}[hint]
Ordering --- volatile stops compiler reordering, not CPU
reordering; you need release/acquire.
\end{solution}
\begin{exercise}[medium]
Estimate the cache-line cost of head and tail sharing one line
(false sharing) on our ring, and name the fix.
\end{exercise}
\begin{solution}[hint]
Each push and pop bounce the line between cores --- pad head and
tail to separate lines.
\end{solution}
\begin{exercise}[hard]
Design a low-latency request path with a 1\,ms p99 budget, from
accept to reply. Budget each hop, choose the data structures,
and defend the two hardest trade-offs.
\end{exercise}
\begin{solution}[hint]
Pools + SPSC rings + pre-allocated reply buffers; the hard
trades are capacity vs latency and batching vs immediacy.
\end{solution}
